\section{Functions: \texttt{lambda} and Application}
\label{sec:lambda}

\subsection{Syntax}

The last part of the lecture adds functions:
\begin{center}
\begin{tabular}{@{}lll@{}}
\toprule
syntax & meaning & abstract syntax \\
\midrule
\code{(lambda x e)} & a function with parameter \code{x} and body \code{e} & \code{lambdaC} \\
\code{(e1 e2)} & call the function \code{e1} on the argument \code{e2} & \code{appC} \\
\bottomrule
\end{tabular}
\end{center}
\begin{lstlisting}
(define-type ArithE
  [numC (n : number)]
  [idC (x : symbol)]
  [plusC (e1 : ArithE) (e2 : ArithE)]
  [timesC (e1 : ArithE) (e2 : ArithE)]
  [letC (x : symbol) (e1 : ArithE) (e2 : ArithE)]
  [lambdaC (x : symbol) (e : ArithE)]
  [appC (e1 : ArithE) (e2 : ArithE)])
\end{lstlisting}
\code{lambdaC} \emph{defines} a function and \code{appC} \emph{uses} one, much as \code{letC} introduces a name and \code{idC} uses it. Functions take one argument; several arguments can be simulated by currying (\Cref{sec:hof}).

\subsubsection{Parsing} Any list whose head is not \code{+}, \code{*}, \code{let}, or \code{lambda} is an application:
\begin{lstlisting}
(define (parse (s : s-expression)) : ArithE
  (cond
    [(s-exp-number? s) (numC (s-exp->number s))]
    [(s-exp-symbol? s) (idC (s-exp->symbol s))]
    [(s-exp-list? s)
     (let [(l (s-exp->list s))]
       (if (s-exp-symbol? (first l))
           (case (s-exp->symbol (first l))
             [(+) (plusC (parse (second l)) (parse (third l)))]
             [(*) (timesC (parse (second l)) (parse (third l)))]
             [(let) (letC (s-exp->symbol (second l)) (parse (third l)) (parse (fourth l)))]
             [(lambda) (lambdaC (s-exp->symbol (second l)) (parse (third l)))]
             [else (appC (parse (first l)) (parse (second l)))])
           ;; head is not a symbol, e.g. ((lambda x (+ x 1)) 5)
           (appC (parse (first l)) (parse (second l)))))]))
\end{lstlisting}
\begin{lstlisting}
(parse '((lambda x (+ x 1)) 5))
;; = (appC (lambdaC 'x (plusC (idC 'x) (numC 1))) (numC 5))
(parse '(f 5))
;; = (appC (idC 'f) (numC 5))
\end{lstlisting}
The outer \code{if} matters for calls whose head is itself a list, like \code{((lambda x (+ x 1)) 5)}: calling \code{s-exp->symbol} on it would fail with \code{s-exp->symbol: not a symbol}.

\begin{remark}
With an application fallback, \code{'(- 1 2)} no longer fails to parse. The \code{else} branch reads it as a call of a function named \code{-}, giving \rkt|(appC (idC '-) (numC 1))| and silently dropping the \code{2}. Before \code{lambda} was added, it was \code{case: no matching clause}.
\end{remark}

\subsection{\texttt{lambda} in Substitution}

A \code{lambda} binds its parameter in its body, exactly as \code{let} binds \code{x} in \code{e2}. So the \code{lambdaC} case of each substitution function is its \code{letC} case without an \code{e1}: stop at shadowing, and in \code{subst-lazy-v2} rename with \code{gensym} first. The application case just recurses.
\begin{lstlisting}
;; in subst-eager (subst-lazy is the same with subst-lazy)
[lambdaC (y e)
         (if (equal? x y) (lambdaC y e) (lambdaC y (subst-eager v x e)))]
[appC (e1 e2) (appC (subst-eager v x e1) (subst-eager v x e2))]

;; in subst-lazy-v2
[lambdaC (y e)
         (if (equal? x y)
             (lambdaC y e)
             (let [(z (gensym))]
               (lambdaC z (subst-lazy-v2 v x (subst-lazy (idC z) y e)))))]
[appC (e1 e2) (appC (subst-lazy-v2 v x e1) (subst-lazy-v2 v x e2))]
\end{lstlisting}

\begin{pitfall}[The pattern variable shadows the argument]
In \rkt|[lambdaC (y e) ...]|, the pattern variable \code{e} (the body) shadows \code{subst}'s own argument \code{e} (the whole \code{lambda}). So the shadowing branch must rebuild \rkt|(lambdaC y e)|; returning plain \code{e} would drop the \code{lambda} and return only its body. Shadowing bites the interpreter's author too.
\end{pitfall}

\subsection{A New Result Type}

A program can now produce a function, for example \code{(lambda x (+ x 1))} or \code{(let f (lambda x x) f)}, and a function is not a number. So evaluation returns a \tyname{Value}, which is either:
\begin{lstlisting}
(define-type Value
  [numV (n : number)]
  [funV (x : symbol) (e : ArithE)])
\end{lstlisting}
\code{+} and \code{*} only make sense on numbers, so they unwrap \code{numV} first and complain about anything else:
\begin{lstlisting}
(define (num+ (v1 : Value) (v2 : Value)) : Value
  (if (and (numV? v1) (numV? v2))
      (numV (+ (numV-n v1) (numV-n v2)))
      (error 'eval "not a number")))

(define (num* (v1 : Value) (v2 : Value)) : Value
  (if (and (numV? v1) (numV? v2))
      (numV (* (numV-n v1) (numV-n v2)))
      (error 'eval "not a number")))
\end{lstlisting}
This is the first \emph{dynamic} type error in our language: whether \code{(+ e1 e2)} is meaningful depends on what \code{e1} and \code{e2} evaluate to, which the interpreter only learns at run time.

\subsection{Eager Application}
\label{sec:eagerapp}

The lecture first adds functions to the \emph{eager} interpreter (\code{lec4b.rkt}, \code{lec4c.rkt}). Application is \emph{call-by-value}: evaluate the function, evaluate the argument, then substitute the argument's \emph{value} into the body. \code{let} works the same way.

\code{subst-eager} substitutes only numbers, but a value may now be a function. So the value is turned back into an expression, and substituted with the expression-substituting \code{subst-lazy} or \code{subst-lazy-v2}. The expression they receive is already evaluated, so nothing is lazy about it:
\begin{lstlisting}
(define (value->expr (v : Value)) : ArithE
  (type-case Value v
    [numV (n) (numC n)]
    [funV (x e) (lambdaC x e)]))

(define (eval-eager-v2 (e : ArithE)) : Value
  (type-case ArithE e
    [numC (n) (numV n)]
    [idC (x) (error 'eval "unbound variable")]
    [plusC (e1 e2) (num+ (eval-eager-v2 e1) (eval-eager-v2 e2))]
    [timesC (e1 e2) (num* (eval-eager-v2 e1) (eval-eager-v2 e2))]
    [letC (y e1 e2)
          (eval-eager-v2 (subst-lazy (value->expr (eval-eager-v2 e1)) y e2))]
    [lambdaC (x e) (funV x e)]
    [appC (e1 e2)
          (type-case Value (eval-eager-v2 e1)
            [funV (x e)
                  (eval-eager-v2 (subst-lazy (value->expr (eval-eager-v2 e2)) x e))]
            [numV (n) (error 'eval "not a function")])]))
\end{lstlisting}
The lecture's \code{lec4b.rkt} writes this with \code{Expr} for \tyname{ArithE} and \code{lambdaV} for \code{funV}, and gives \code{subst} a \tyname{Value} that it converts at the \code{idC} case; the effect is the same. Functions can now be passed around:
\begin{lstlisting}[style=repl]
(let do-twice (lambda f (f (f 1)))
  (let double (lambda x (* 2 x))
    (do-twice double)))
---> (do-twice double)  with both substituted
---> (double (double 1))
---> 4
\end{lstlisting}

\subsubsection{Eager capture} \Cref{sec:lazy} said that eager substitution cannot capture because a number has no names in it. A function value does have names, and its body may contain free variables:
\begin{lstlisting}[style=repl]
(let f (lambda x (+ x y))         y is free in f's body
  (let y 6
    (f 1)))
---> (let y 6 ((lambda x (+ x y)) 1))   substitute f's value: y is captured
---> ((lambda x (+ x 6)) 1)
---> 7
\end{lstlisting}
The right answer is an error: the \code{y} in \code{f} refers to no binding, and Racket rejects the same program with \code{y: unbound identifier}. \code{eval-eager-v2} returns 7 because \code{subst-lazy} carried the free \code{y} under \code{(let y 6 ...)}. Eager evaluation is not a defense against capture once values can contain code.

The fix (\code{lec4c.rkt}) is the one from \Cref{sec:capavoid}: \code{eval-eager-v3} is \code{eval-eager-v2} with both calls to \code{subst-lazy} replaced by \code{subst-lazy-v2}. Now the program above raises \code{unbound variable}, and returning \code{f} instead of calling it shows the renaming:
\begin{lstlisting}[style=repl]
> (eval-eager-v3 (parse '(let f (lambda x (+ x y)) (let y 6 f))))
(funV 'g3507 (plusC (idC 'g3507) (idC 'y)))
\end{lstlisting}
Substituting \code{f} into the body of \code{(let y 6 ...)} renamed \code{f}'s parameter to a fresh name, as the \code{lambdaC} case of \code{subst-lazy-v2} always does; the free \code{y} stays free.

\subsection{Lazy Application: \texttt{eval-lazy-v2}}

The last version (\code{lec4d.rkt}) makes application lazy, like \code{let} in \Cref{sec:lazy}: the argument is substituted unevaluated.

\begin{lstlisting}
(define (eval-lazy-v2 (e : ArithE)) : Value
  (type-case ArithE e
    [numC (n) (numV n)]
    [idC (x) (error 'eval "unbound variable")]
    [plusC (e1 e2) (num+ (eval-lazy-v2 e1) (eval-lazy-v2 e2))]
    [timesC (e1 e2) (num* (eval-lazy-v2 e1) (eval-lazy-v2 e2))]
    [letC (y e1 e2) (eval-lazy-v2 (subst-lazy-v2 e1 y e2))]
    [lambdaC (x e) (funV x e)]
    [appC (e1 e2)
          (type-case Value (eval-lazy-v2 e1)
            [funV (x e) (eval-lazy-v2 (subst-lazy-v2 e2 x e))]
            [numV (n) (error 'eval "not a function")])]))
\end{lstlisting}
\begin{itemize}
  \item \textbf{\code{lambdaC}}: a \code{lambda} is already a value. There is nothing to compute until it is called.
  \item \textbf{\code{appC}}: evaluate \code{e1}; it must be a function \rkt|(funV x e)|. Then call it the way \code{let} works: substitute the argument \code{e2} for \code{x} in the body \code{e}, and evaluate the result. \code{e2} is substituted unevaluated, because this is the lazy version.
\end{itemize}
\begin{lstlisting}[style=repl]
((lambda x (+ x 1)) 5)
---> (+ 5 1)             subst 5 for x in the body
---> 6
\end{lstlisting}

\begin{keypoint}[\code{let} is application]
\code{((lambda x e) e2)} means the same as \code{(let x e2 e)}. What \code{lambda} adds is that the function is a \emph{value}: it can be named with \code{let}, passed as an argument, and returned as a result.
\end{keypoint}

\begin{center}
\begin{tabular}{@{}ll@{}}
\toprule
program & result \\
\midrule
\code{(let f (lambda x (* x 2)) (f 5))} & \code{(numV 10)} \\
\code{((lambda f (f 3)) (lambda x (+ x 1)))} & \code{(numV 4)}: a function as argument \\
\code{(((lambda x (lambda y (+ x y))) 1) 2)} & \code{(numV 3)}: a function as result \\
\code{((lambda x 5) y)} & \code{(numV 5)}: the unused argument is never evaluated \\
the \code{do-twice} example of \Cref{sec:eagerapp} & \code{(numV 4)} \\
\code{(5 3)} & \code{not a function} \\
\code{(+ (lambda x x) 1)} & \code{not a number} \\
\bottomrule
\end{tabular}
\end{center}

\subsection{Capture, Again}

With functions, capture-avoiding substitution matters even more. A function can return a function, so a free name can end up under another \code{lambda}:
\begin{lstlisting}[style=repl]
(((lambda x (lambda y (+ x y))) y) 2)
subst-lazy:     (lambda y (+ y y)),  then 2 ---> 4          captured!
subst-lazy-v2:  (lambda g (+ y g)),  then 2 ---> (+ y 2) ---> unbound variable
\end{lstlisting}
\code{eval-eager} and \code{eval-lazy} still return a \code{number}, and \code{subst-eager} can substitute only numbers, so their \code{lambdaC} and \code{appC} cases remain errors (\code{"TODO: lambda"}); \code{eval-eager-v2}/\code{v3} and \code{eval-lazy-v2} are the versions with functions.

\begin{center}
\begin{tabular}{@{}lll@{}}
\toprule
program & \code{eval-eager-v3} (call-by-value) & \code{eval-lazy-v2} (call-by-name) \\
\midrule
\code{((lambda x 5) y)} & \code{unbound variable} & \code{(numV 5)} \\
\code{(let x (+ y y) 10)} & \code{unbound variable} & \code{(numV 10)} \\
the \code{f}/\code{y} example of \Cref{sec:eagerapp} & \code{unbound variable} & \code{unbound variable} \\
\bottomrule
\end{tabular}
\end{center}
With capture-avoiding substitution, the two strategies differ only where an argument is never used, as for \code{let} in \Cref{sec:lazy}. Racket, like most languages, is call-by-value.
